\scanpage{22}
\mspage{\textcircled{15}}
On fixe entier $n$ premier aux car. résiduelles

\uline{Classe fondamentale locale}

Supposons $Y\hookrightarrow X$ ``intersection complète de dim relative $1$'' et supposons d'abord $Y=V(\varphi)$, $\varphi$ section régulière de $\mathcal O_X$. Alors si $\unsure{u|U}:X-Y=U\longrightarrow X$, donc \deleted{$i_*(\underline O_U)\sim O_X[\frac{1}{\varphi}]$}, \deleted{$i_*(\underline O^*_U)$ et} $\varphi\in\Gamma(i_*(\deleted{\unsure{\underline O}}\underline G_{mU}))$ \deleted{et $\unsure{\sim|=|\to}\ \in\Gamma(\unsure{X;|i_*}\,i_*(\unsure{G|O}^*_{mU}))$} \deleted{D'\unsure{où|ai|au}} ailleurs $\varphi$ est \unsure{déterminé|déterminée} \inserted{par $Y$} mod multiplication par un $\in\Gamma(X,i_*(\deleted{\unsure{\underline O}}G_{mX}))$, donc \deleted{l'image} $\varphi$ définit une section \deleted{$\varphi\in$} \uline{canonique} de \deleted{$i_*(\underline O^*_U)/$} $i_*(G_{mU})/G_{mX}$. Or on a la suite exacte canonique

\diagram{\begin{tikzcd}[ampersand replacement=\&,column sep=.75em,row sep=.5em]
0 \arrow[r] \& \underline H^0_Y(G_{mX}) \arrow[r]\arrow[d,equal] \& G_{mX} \arrow[r,dash] \& i_*(G_{mU}) \arrow[r,dash,"d"] \& \underline H^1_Y(G_{mX}) \arrow[d] \\
\& 0 \& \& \& 0
\end{tikzcd}}

qui donne ici

\[
\deleted{\unsure{H^1|H}}\ \underline H^1_Y(G_{mX})\xrightarrow{\sim}i_*(G_{mU})/\mathrm{Im}\,G_{mX}\quad \unsure{:|;})
\]

et montre que $d\varphi$ est une section \uline{canonique} de $\underline H^1_Y(G_{mX})$, que nous allons noter $\deleted{\unsure{\xi|\gamma|\varepsilon}_{Y/X}}\underline\beta_{Y/X}$. D'autre part, on a \deleted{une}\inserted{la} suite exacte de Kummer

\[
\unsure{0|\cdot}\mathrel{-}\unsure{\mu_{n_X}|\mu_{nX}}\longrightarrow G_{mX}\overset{n}{\mathrel{-}} G_{mX}\longrightarrow 0
\]

% archive sheet number: 21
